% Generated by evidence/summarize_ara_v1.py; do not edit.
\begin{table*}[t]
\centering\footnotesize
\caption{本地 point-v1 归档协议与六维搜索分布。}
\label{tab:ara-v1-protocol}
\renewcommand{\arraystretch}{1.12}
\begin{tabularx}{\textwidth}{@{}lYY@{}}
\toprule
项目 & 设置 & 证据边界 \\
\midrule
模型 & Qwen3.8-27B（本地路径） & 模型提交未记录；64 层 \\
硬件 & RTX PRO 6000 Blackwell Server Edition & 97,887 MiB；driver 595.58.03 \\
加载 & BF16；无量化；自动放置 & GPU 配置上限 90 GiB \\
无害数据 & mlabonne/harmless\_alpaca & revision: 02c6a92cfcf1（完整标识见证据 JSON） \\
有害数据 & mlabonne/harmful\_behaviors & revision: 01cead013989（完整标识见证据 JSON） \\
校准 & 每类候选 train[:300]；实际 64+64 & 捕获批量 1；秩至多 128 \\
验证 & 每类 train[300:400]；各 100 条 & 120 次搜索重复使用 \\
审计 & 每类 test[:100] & 仅配置；无审计得分 \\
解码 & seed=42；不采样；关闭 thinking & 最多 100 词元；有效批量 16 \\
批量配置 & 自动配置 0；最大值 16 & 有效评估批量为 16 \\
优化 & L-BFGS：一次 step；max\_iter=20 & history=10；strong Wolfe \\
采样 & TPE：36 startup + 84 adaptive & 候选 128；multivariate；seed=42 \\
数值约束 & 全部 constraints=[0.0] & 仅无记录数值失败；不代表效果验收 \\
起始层比例 & [0.25, 0.65] & 线性采样 \\
层跨度比例 & [0.1, 0.55] & 线性采样 \\
注意力强度 & [0.001, 1.0] & 对数采样 \\
MLP 原始强度 & [-0.1, 0.5] & 线性采样 \\
推离权重 & [0.0, 2.0] & 线性采样 \\
间隔 & [0.25, 4.0] & 对数采样 \\
区间解析 & floor(64×起点)，ceil(64×跨度) & 结束层裁剪至 64；MLP 强度取 max(0,原值) \\
\bottomrule
\end{tabularx}
\end{table*}
